\chapter{Introduzione alla bioinformatica}\label{chapter:cap_one}
Dall'alba dei tempi tutto cio che circonda l'uomo è stato fonte di curiosità e studio.
Partendo dall'interesse di tutti i vari aspetti riguardanti l’essere umano e gli organismi viventi,
con cui ci relazioniamo quotidianamente, si è sviluppata la branca della biologia

L'informatica, invece, il branco della scienza che si occupa dello studio e dell'elaborazione delle informazioni,
è stata fin da sempre utilizzata come strumento per migliorare la qualità della vita dell'uomo\cite{treccaniInformatica}.
Dalla prima calcolatrice meccanica, la \textit{Pascalina}, inventata da Blaise Pascal nel 1644\cite{pascal1995pensieri}, 
alla \textit{Macchina di Turing}, ideata da Alan Turing nel 1943\cite{singh1999codici}, 
fino ai moderni computer quantistici, l'informatica ha sempre avuto un ruolo fondamentale nell'evoluzione della società umana.

Lo sviluppo di entrambe queste discipline col passare degli anni ha portato ad un'innovazione: la \textit{bioinformatica}.
La bioinformatica nasce dalla combinazione delle due materie prima elencate in quanto gli scienziati avevano
sempre più l'esigenza di gestire i dati rilevati dalla biologia
molecolare, spesso file di testo che superano le decine di gigabyte, in modo da poter interpretare più facilmente i campioni raccolti.
\section{Gene}\label{sec:cap_one_bioinfo}
Ogni organismo eucariote o procariote pluricellulare o unicellulare che sia presenta una molecola
fondamentale che contiene tutte le informazioni necessarie per costituire l'organismo, tale molecola
è il DNA (acido desossiribonucleico).
Il DNA è composto da subunità, chiamate nucleotidi\footnote{Il nucleotide è l'unità di base che forma DNA e RNA}, contenenti: un gruppo fosfato, uno zucchero
pentoso e una base azotata.
Una specifica sequenza di DNA forma un gene e, l'insieme di tutti i geni, genera il genoma che, oltre
a racchiudere le informazioni genetiche, presenta anche le parti del DNA che non codificano proteine\footnote{Le proteine sono macromolecole che si formano
partendo dal DNA e che vanno a costituire i tessuti.}.
Circa il 30\% del DNA nel genoma umano si trova negli esoni e negli introni. I primi, gli esoni, sono
sequenze che codificano gli aminoacidi, i quali, vengono espressi in proteine, i secondi, gli introni,
invece non codificano nessuna proteina. Al giorno d'oggi sappiamo che il genoma umano presenta
circa 20000 geni che codificano proteine.
Con il passare degli anni sono stati sequenziati molti genomi, sia di microrganismi che di animali,
questo perché, l'interesse degli scienziati è quello di andare definire i processi evolutivi degli esseri
viventi, ma, per ottenere dei risultati, è indispensabile andare a comparare numerosi genomi di
organismi differenti.

La branca della bioinformatica è molto ampia, per questo, noi ci interesseremo
principalmente del genoma.
La bioinformatica ha un ruolo
fondamentale negli studi della struttura del genoma, per esempio, nell'assemblaggio 
e annotazione di genomi, o nella caratterizzazione del profilo di espressione genica,
generando informazioni cruciali per orientare i successivi studi di genomica funzionale che si avvalgono di una grande varietà di tecniche sperimentali\cite{fondamentiBioinformatica}

Tutte le cellule di un organismo hanno sostanzialmente lo stesso genoma
ma possono avere epigenomi e trascrittomi molto diversi.
\subsection{Epigenoma}\label{subsec:epigenomi}
L'\textit{epigenoma} è una moltitudine di composti chimici che che regolano il genoma su cosa fare senza cambiare la sequenza del DNA.
Esso ha le varie istruzioni per poter costuire le proteine che servono per svolgere le diverse funzioni della cellula.
L'epigenoma poi, costituito da composti chimici e proteine, può attaccarsi al DNA e impartire diversi comandi, tra cui l'attivazione o la disattivazione di geni, oppure controllare
la produzione di determinate proteine dentro la cellula.

Un essere umano ha trilioni di cellule, specializzate per diverse funzioni, nei muscoli, 
nelle ossa e nel cervello, e ciascuna di queste cellule ha essenzialmente lo stesso genoma nel suo nucleo. 
Le differenze tra le cellule sono determinate da come e quando diversi gruppi di geni vengono attivati 
o disattivati in vari tipi di cellule. Le cellule specializzate nell'occhio attivano geni che producono 
proteine in grado di rilevare la luce, mentre le cellule specializzate nei 
globuli rossi producono proteine che trasportano l'ossigeno dall'aria al resto del corpo. 
L'epigenoma controlla molti di questi cambiamenti nel genoma.\cite{genomeGov}

\subsection{Trascrittoma}\label{subsec:trascrittomi}
Il DNA, non viene utilizzato direttamente per costruire le proteine: prima deve essere trascritto in RNA. 
Questo processo prende il nome di trascrizione e produce diversi tipi di RNA.
Il \textit{trascrittoma} rappresenta l'insieme di tutte le molecole di RNA presenti in una cellula in un dato momento.

L'RNA principale è RNA messaggero (mRNA). che trasporta le istruzioni dai geni ai ribosomi, 
le “fabbriche” molecolari della cellula. Nei ribosomi, 
la sequenza di basi dell'mRNA viene letta e tradotta in catene di amminoacidi, 
i mattoni che formano le proteine.
Accanto all'mRNA, esistono anche altri tipi di RNA che non codificano proteine ma che svolgono funzioni regolatorie 
o strutturali, contribuendo a controllare l'attività dei geni e l'organizzazione della cellula.
Poiché ogni cellula del corpo umano possiede lo stesso patrimonio genetico, 
ciò che la distingue dalle altre non è il genoma ma l'insieme dei geni che vengono attivati o silenziati. 
Questo insieme di trascrizioni geniche costituisce il trascrittoma, e varia notevolmente da cellula a cellula 
e da tessuto a tessuto. Ad esempio, le cellule muscolari esprimono geni che producono proteine contrattili, 
mentre le cellule nervose attivano geni che permettono la trasmissione dei segnali elettrici.

Il trascrittoma quindi è lo specchio dell'attività del genoma: mostra quali geni sono effettivamente 
in funzione in un determinato momento e fornisce informazioni preziose sull'identità, 
lo stato e la salute delle cellule.\cite{genomeGov}

\section{Trascrittomica spaziale}\label{sec:cap_one_bioinfo}
Come detto in precendenza, il trascrittoma rappresenta l'insieme di tutte le molecole di RNA presenti in una cellula in un dato momento.
La trascrittomica a singola cellula (scRNA-seq) è diventata essenziale per la ricerca biomedica nell'ultimo decennio, 
in particolare nello sviluppo della biologia, cancro, immunologia e neuroscienze. 

La maggior parte dei protocolli scRNA-seq disponibili in commercio richiede che le cellule siano recuperate intatte e vitali dai tessuti. 
Ciò ha precluso lo studio di molti tipi cellulari e ha in gran parte distrutto il contesto spaziale che altrimenti avrebbe potuto 
informare le analisi dell'identità e della funzione cellulare. 
Oggi, un numero crescente di piattaforme disponibili in commercio facilita la valutazione altamente dimensionale e spazialmente 
risolta della trascrizione genica, nota come \textit{trascrittomica spaziale (TS)}.\cite{trascrittomicaSpazialeIntroduzione}

Una sfida fondamentale nella TS è l'identificazione degli \textit{Spatially Variable Genes (SVGs)}, e degli \textit{Highly Variable Genes (HVGs)}.
